\documentclass[11pt]{article}
\usepackage[a4paper,margin=25mm]{geometry}
\usepackage{amsmath,amssymb,amsthm}
\usepackage{booktabs}
\usepackage{array}
\usepackage{setspace}
\usepackage[tablesonly,nolists]{endfloat}
\usepackage[round]{natbib}
\usepackage[T1]{fontenc}
\usepackage{lmodern}
\usepackage{url}
\newcommand{\doi}[1]{\url{https://doi.org/#1}}
\usepackage[hidelinks]{hyperref}
\hypersetup{pdftitle={Private Innovation Intermediation, Regional Incidence, and Public Access},pdfauthor={},pdfsubject={Anonymous review manuscript}}
\title{Private Innovation Intermediation,\\ Regional Incidence, and Public Access}
\author{}
\date{}
\begin{document}
\maketitle
\begin{abstract}
How do private innovation fees and local public-access decisions depend on regional commercialisation returns, service values and technical complementarity, and when do local and national provision evaluations diverge? We study two local governments and two private intermediaries with endogenous fees, founder selection, cross-region research matching and differentiated Cournot competition. Private pricing selects abilities and free basic access reduces private receipts without necessarily inducing exit. Conditional on full participation in the partner region, we derive separate local and national public-cost ceilings. A primitive profit-share condition shows that stronger technical complementarity can raise the local ceiling even when additional ventures reduce regional commercialisation surplus. This result is analytical for sufficiently weak matching; exact rational certification establishes its signs on a continuous two-parameter region at the original finite intensity. Three separately certified four-party entry examples compare decentralised provision with national welfare maximisation over sixteen configurations: two exhibit public underprovision and one coincidence. Their configuration switch requires a narrow public-cost band and fails a recorded alternative ability distribution. The paper provides conditional, model-specific policy boundaries, separating familiar screening and externality mechanisms from the candidate contribution of their regional-incidence interaction.
\end{abstract}
\noindent\textbf{Keywords:} Innovation intermediation; Public access; Participant selection; Mixed provision; Regional welfare.

\noindent\textbf{JEL classification:} D43; H77; L31; O32.
\doublespacing

\input{manuscript/sections/01_introduction}
\input{manuscript/sections/02_model}
\input{manuscript/sections/03_continuation}
\input{manuscript/sections/04_propositions}
\input{manuscript/sections/05_numerical}
\input{manuscript/sections/04_diagnostic}
\input{manuscript/sections/06_empirical_implications}
\input{manuscript/sections/06_related_literature}
\input{manuscript/sections/07_institutions}
\input{manuscript/sections/09_conclusion}
\input{manuscript/sections/08_declarations}
\begingroup\raggedright\interlinepenalty=10000
\bibliographystyle{plainnat}
\bibliography{manuscript/references}
\endgroup
\end{document}
